\documentstyle{article}
\begin{document}

\centerline{\bf Concepts and Algorithms of Computer Analysis}
\centerline{\bf for an Approximate Solution of ODE's}

\vskip 0.3cm

\centerline{Karl Hantzschmann}
\centerline{Universitaet Rostock, Fachbereich Informatik}
\centerline{Albert-Einstein-Strasse 21}
\centerline{18051 Rostock, Germany}
\centerline{Email: hantzschmann@informatik.uni-rostock.de}

\vskip 0.3cm

1.	Purpose and methodology of Computer Analysis

Whereas in Computer Algebra the dominant object is algorithms based on
algebraic concepts and aimed  at finding exact solutions, Computer Analysis
gives priority to algorithms using CAS for finding controlled analytic
approximate solutions. The point is to find efficiently practicable and
transparent formula expressions reflecting the qualitative properties of a
given problem and securing a desired accuracy. This implies practicable
error estimations which must always be performed by the Computer.

2.	Algorithms for approximate solutions

In the field of ODEs several algorithms could be developed and tried in the
last decade. In accordance with the purpose of Computer Analysis preference
is given to mathematical approximation methods allowing an automatic
adaptation to a given problem and giving a transparent short formula
expression at a reasonable expense. The paper will give a survay of some
these algorithms.
In recent years good results could be achieved by a two-step concept. In the
first step the given problem will be adapted by means of a neighbouring
problem which must be suitably chosen and allow to find an exact solution by
algorithms of Computer Algebra. Then the fundamental solutions form the
basis for an adequate approximating ansatz in the second step. This two-step
concept has the advantage of allowing a close problem adaptation. It offers
a wide field for the development of several algorithms depending first on
the approximation criterions used in both steps and second on the selected
family of neighbouring problems. Interpreting the steps as projection
methods allows a uniform theoretic foundation.

A modification of this strategy results in an algorithm for the approximate
computation of non-exponential solutions for the class of simple
differential equations.
Based on efficient algorithms for the finding exponential solutions or some
subclasses of such solutions a relatively wide class of adapted differential
equations has been made available which may serve as a basis for
approximations. This class consists of simple differential equations, the
relevant differential operators of which are factorizable by differential
operators of first order with exponential solutions.

3.	Practicable error estimations

For finding error estimations a special procedure has proved its worth.
This procedure will allow to completely calculate error boundaries by
computer, if the coefficient functions of the differential equations belong
to a reasonably restricted function class.  It will be necessary to
formulate the error function in the form of a fix point equation which will
use Green's function of a linear neighbouring differential equation which
in its turn can be derived from the given differential operator by means of
linearization.

If the non-linear remainder operators satisfy special Lipschitz-relations
of second order in surroundings of the approximate solution, fixpoint
theorems will give the wanted estimations. In the meantime the principle has
been modified and extended. It will be discussed in the paper by taking a
non-linear differential equation of n-th order as an example.
This procedure shows the advantage of combining the symbolic with the proven
numerical computation as a typical attribute of the algorithms of Computer
Analysis.
The Lipschitz-estimations needed for non-linear differential equations are
supplied by a specially developed Lipschitz-calculus.

\end{document}
